%%%%%%%%%%%%%%%%%%%%%%%%%%%%%%%%%%%%%%%%%
% Medium Length Professional CV
% LaTeX Template
% Version 2.0 (8/5/13)
%
% This template has been downloaded from:
% http://www.LaTeXTemplates.com
%
% Original author:
% Trey Hunner (http://www.treyhunner.com/)
%
%%%%%%%%%%%%%%%%%%%%%%%%%%%%%%%%%%%%%%%%%

\documentclass{resume}
\usepackage[left=0.75in,top=0.6in,right=0.75in,bottom=0.6in]{geometry}
\usepackage{xeCJK}
\usepackage{color}
\usepackage[UTF8]{ctex}
\usepackage[citecolor=green]{hyperref}

\name{陈柏润}
\address{+(86) 15217720601 \\ chenborun1101@gmail.com \\ github.com/PorunC}

\begin{document}

%----------------------------------------------------------------------------------------
%	PROFESSIONAL SUMMARY
%----------------------------------------------------------------------------------------

\begin{rSection}{个人概述}

3年后端开发经验，专注将大模型落地为生产级业务工具。
主导过AI周报系统及代码智能平台三个从0到1的项目，具备独立完成系统设计、
全链路开发及持续交付的能力。

\end{rSection}

%----------------------------------------------------------------------------------------
%	WORK EXPERIENCE SECTION
%----------------------------------------------------------------------------------------

\begin{rSection}{工作经历}

\begin{rSubsection}{招银网络科技有限公司}{2023.7 -- 至今}
{后端开发工程师 | 核心应用开发团队}

\item \textbf{AI 周报系统（独立负责）} \\
【情境】银行运维数据周报依赖人工从多源系统采集汇总，每轮需8-10小时，效率低且易出错。\\
【任务】构建全自动化的LLM驱动周报系统。\\
【行动】从0自研多源异构数据采集、LLM智能分段分析、周报组装与推送的全链路闭环，内置重试容错与事务保障。\\
【结果】周报耗时从8-10小时降至20分钟，年节省50人天，投产至今稳定运行。

\item \textbf{CodeWiki 代码智能平台（独立负责）} \\
【情境】跨模块代码理解依赖人工逐行阅读源码，新人上手需数周，研发效率受影响。\\
【任务】建设代码理解与知识管理平台，降低团队代码认知成本。\\
【行动】自研多种编程语言AST解析引擎与多维依赖图谱，基于GraphRAG实现语义级代码检索，驱动LLM自动生成源码级文档；全栈自研（Python/FastAPI + TypeScript/React），支持多种交互。\\
【结果】日均查询50+次，跨模块理解时间从2小时缩至30分钟，覆盖日常代码评审与新人onboarding。

\item \textbf{AI Agent 长期记忆与上下文压缩系统（独立负责）} \\
【情境】AI Agent 在长任务和多轮对话中缺乏稳定的长期记忆能力，历史上下文与工具日志持续膨胀，导致推理准确率下降、Token 成本上升。\\
【任务】设计并落地一套生产级 Agent 记忆系统，提升长期任务的上下文利用率、记忆召回准确率与执行稳定性。\\
【行动】从0设计 L0-L3 多层记忆架构，将原始对话、结构化摘要、语义记忆与长期知识分层管理；构建基于向量检索、BM25 全文检索与 Hybrid RRF 融合召回的记忆检索链路，支持 SQLite 本地存储与云端向量数据库双后端；设计上下文卸载与工具日志压缩机制，对长任务执行过程进行结构化摘要、追踪与恢复；配套自动调度、会话清理、指标上报和毫秒级性能计时，保障系统可观测与可维护。\\
【结果】系统已投产并覆盖多个业务渠道，Token 消耗降低 61\%，任务通过率相对提升 52\%，长期记忆准确率由 48\% 提升至 76\%。
\end{rSubsection}

\end{rSection}

%----------------------------------------------------------------------------------------
%	SKILLS SECTION
%----------------------------------------------------------------------------------------


%----------------------------------------------------------------------------------------
%	EDUCATION SECTION
%----------------------------------------------------------------------------------------

\begin{rSection}{教育}

{\textbf{深圳大学}} \hfill {\em 2019.9 -- 2023.6} \\
本科生，软件工程

\end{rSection}

%----------------------------------------------------------------------------------------
%	OTHERS SECTION
%----------------------------------------------------------------------------------------

\begin{rSection}{其他}

\begin{rSubsection}{}{}{}{}
\item[-] 英语：CET-6，可阅读及编写英文技术文档
\item[-] 博客：misaka-9982.com
\end{rSubsection}

\end{rSection}

\end{document}
